% ======================================================================
% SECTION 3: AUTONOMOUS TRIAL DESIGN AND STATISTICAL FRAMEWORK
% File: sections/trial_design.tex
% ======================================================================

[PROCESSING INSTRUCTION: Generate this section (approximately 3-4 pages) covering autonomous
trial design, protocol generation, and statistical planning. Use the following source files:

SOURCE FILES TO PROCESS:
- sponsor/input\_files/sponsor\_03\_trial\_design\_stats.md (6 design patterns, biomarker strategy,
  adaptive/Bayesian statistics, estimand framework)
- sponsor/input\_files/org\_03\_functions\_discovery\_protocol.md (IND-enabling, protocol design,
  ICH M11 structured protocol, adaptive design governance)
- sponsor/input\_files/sponsor\_02\_preclinical\_regulatory.md (translational package, IND mechanics)
- national-platform/new\_paper/final\_paper/sections/ich\_e6r3\_adaptation.tex (adapted GCP)
- national-platform/new\_paper/final\_paper/sections/cfr312\_adaptation.tex (adapted IND process)

SUBSECTION STRUCTURE:

\subsection{Clinical Accountability Agent (clinical\_accountability\_agent)}
- Protocol medical logic automation, eligibility rule generation
- Dose/schedule logic using Project Optimus requirements
- Medical monitoring rule definition
- Maps to: Study Responsible Physician / Clinical Lead

\subsection{Autonomous Protocol Design}
- ICH M11 structured protocol generation
- 6 design patterns (from sponsor\_03): conventional RCT, seamless phases,
  basket master protocol, umbrella master protocol, platform trial, single-arm + external control
- Endpoint/estimand definition using ICH E9(R1) framework
- Adaptive design governance with Bayesian methods (FDA Jan 2026 draft guidance)

\subsection{Biomarker and Companion Diagnostics Strategy}
- Biomarker funnel math and CDx co-development planning
- Evidence hierarchy for biomarker qualification
- EU IVDR co-regulation requirements
- ctDNA as biomarker integration (from sponsor\_02)

\subsection{Biostats and Data Agent (biostats\_data\_agent)}
- Statistical Analysis Plan (SAP) generation
- Interim analysis planning with futility/efficacy boundaries
- Bayesian adaptive randomization logic
- Dataset lock procedures and CDISC SDTM/ADaM compliance

PYTHON SCRIPT REQUIREMENTS (to be generated by Claude Code Opus 4.6 in the future):
- File: sponsor/template/scripts/clinical\_accountability\_agent.py
  Purpose: Implement eligibility rule engine, dose-schedule optimizer using
  Project Optimus criteria, and medical monitoring rule generator.
- File: sponsor/template/scripts/protocol\_generator.py
  Purpose: Generate ICH M11-compliant protocol documents. Input: indication,
  phase, design pattern, endpoints. Output: structured protocol sections.
- File: sponsor/template/scripts/biostats\_agent.py
  Purpose: SAP generator with Bayesian adaptive methods. Implement interim
  analysis boundaries, futility stopping rules, and adaptive randomization.
  Reference FDA January 2026 Bayesian draft guidance.
- File: sponsor/template/scripts/biomarker\_strategy\_engine.py
  Purpose: Biomarker qualification scoring, CDx timeline planning, ctDNA
  integration logic. Input: assay data, regulatory requirements. Output:
  biomarker development plan with milestones.
- All scripts must pass ruff lint (line-length 120, E/F/W rules).
- Process with Claude Code Opus 4.6 (1M token context).

TABLE REQUIREMENTS:
- Table 3: Trial Design Pattern Comparison (6 patterns from sponsor\_03)
  Columns: Design Pattern | Best For | Key Advantage | Regulatory Basis | Agent Config
- Table 4: Estimand Framework Components (ICH E9(R1))
  Columns: Component | Definition | Agent Implementation | Example

FORMATTING NOTES:
- Include comparative table spanning full page width using tabularx.
- Reference adapted ICH E6(R3) for quality-by-design protocol principles.]

\subsection{Clinical Accountability Agent}

[PROCESSING INSTRUCTION: Generate 1 page on the clinical\_accountability\_agent. This agent
replaces the Study Responsible Physician / Clinical Lead. It automates protocol medical
logic, eligibility criteria generation (inclusion/exclusion based on indication, biomarkers,
prior therapies), dose/schedule optimization per Project Optimus, and medical monitoring
rules. Reference the protocol design/scientific leadership function from
org\_03\_functions\_discovery\_protocol.md. Include the endpoint/estimand definition process
and adaptive design governance. Process with Claude Code Opus 4.6 (1M token context).]

\subsection{Autonomous Protocol Design}

[PROCESSING INSTRUCTION: Generate 1-1.5 pages covering the 6 trial design patterns
from sponsor\_03\_trial\_design\_stats.md as a comparative framework. Include a table
(Table 3) comparing: conventional RCT, seamless phases, basket master protocol, umbrella
master protocol, platform trial, and single-arm + external control. For each design,
describe how the autonomous system selects and configures the appropriate pattern
based on indication characteristics, biomarker prevalence, and regulatory pathway.
Reference ICH M11 structured protocol format. Reference ICH E20 adaptive design draft.
Process with Claude Code Opus 4.6 (1M token context).]

\subsection{Biomarker and Companion Diagnostics Strategy}

[PROCESSING INSTRUCTION: Generate 0.75-1 page on the biomarker strategy. Cover the
evidence hierarchy for biomarker qualification programs, CDx co-development planning
with EU IVDR requirements, and ctDNA integration as a monitoring biomarker. Reference
the biomarker funnel math concept from sponsor\_01 and the translational/nonclinical
package from sponsor\_02. Process with Claude Code Opus 4.6 (1M token context).]

\subsection{Biostatistics and Data Agent}

[PROCESSING INSTRUCTION: Generate 1 page on the biostats\_data\_agent. Cover: SAP
generation with Bayesian adaptive methods per FDA January 2026 draft guidance,
interim analysis design with futility/efficacy boundaries, adaptive randomization,
and CDISC compliance (SDTM/ADaM). Reference ICH E9(R1) estimand framework with
a table (Table 4) showing the five components. Process with Claude Code Opus 4.6
(1M token context).]
